% Source: OpenAI2026AreaLaw, 09-amplification.tex, lines 101--139,
% openai/math@adc7f1241b42e322a6451854ab7e4b4c146bf78a.
% This leaf proves the local-channel and physical shell-oscillation steps.
% It does not assert Lemma 10.1, a one-event increment, or clock propagation.
% Public declarations: LocalRootChannels and ChannelShellOscillation.

\section{Local root channels and physical shell oscillations}
\label{sec:al_channel_shell_oscillation}

Use the physical space $\mathcal H=(\C^q)^{\otimes I}$, the finite
spectator $\mathcal K=\C^J$, and the algebras $\mathcal A_K$ of
Section~\ref{sec:qc_site_oscillation}, with $q>0$. The dimension of
$\mathcal K$ may be zero. Let $k$ be a positive contraction on
$\mathcal H$, so $0\le k\le\mathbb 1$.
Write $\widehat A=A\otimes\mathbb 1_{\mathcal K}$ for a physical
operator $A$. All operators $B,C$ on $\mathcal H\otimes\mathcal K$
are arbitrary complex matrices, with the Euclidean operator norm.
For $K\subseteq I$, the expectation $E_K$ retains exactly the sites
in $K$ and averages the sites in $I\setminus K$.

\subsection{Localized square-root channels}

\begin{definition}[Root channels with a spectator]
    \label{def:al_local_root_channel}
    \uses{def:qc_site_expectation}
    \lean{TNLean.PEPS.AreaLaw.spectatorRootChannel,
        TNLean.PEPS.AreaLaw.localRootChannel,
        TNLean.PEPS.AreaLaw.spectatorRootChannel_sub,
        TNLean.PEPS.AreaLaw.localRootChannel_sub}
    \leanok
    Set $k_K=E_K(k)$, and define the positive square roots
    \begin{align}
        G&=(\mathbb 1-k)^{1/2}, & R&=k^{1/2},
        \label{eq:al_shell_full_roots}
    \end{align}
    \begin{align}
        G_K&=(\mathbb 1-k_K)^{1/2}, & R_K&=k_K^{1/2}.
        \label{eq:al_shell_local_roots}
    \end{align}
    The full and localized channels are
    \begin{align}
        \mathcal E(B)&=\widehat G B\widehat G+\widehat R B\widehat R,
        \label{eq:al_shell_full_channel}
    \end{align}
    \begin{align}
        \mathcal E_K(B)&=\widehat G_K B\widehat G_K
            +\widehat R_K B\widehat R_K.
        \label{eq:al_shell_local_channel}
    \end{align}
    In particular, $\mathcal E(B-C)=\mathcal E(B)-\mathcal E(C)$
    and $\mathcal E_K(B-C)=\mathcal E_K(B)-\mathcal E_K(C)$ by
    distributivity. These are the localized channels of
    \cite[Section~10]{OpenAI2026AreaLaw}, extended by the identity on
    the spectator.
\end{definition}

\begin{theorem}[Contraction of the root channels]
    \label{thm:al_root_channel_contraction}
    \uses{def:al_local_root_channel}
    \lean{TNLean.PEPS.AreaLaw.norm_spectatorRootChannel_le,
        TNLean.PEPS.AreaLaw.norm_localRootChannel_le,
        TNLean.PEPS.AreaLaw.norm_localRootChannel_sub_self_le}
    \leanok
    The full and localized channels satisfy
    \begin{align}
        \|\mathcal E(B)\|&\le\|B\|,
        \label{eq:al_shell_full_contraction}
    \end{align}
    \begin{align}
        \|\mathcal E_K(B)\|&\le\|B\|,
        \label{eq:al_shell_local_contraction}
    \end{align}
    \begin{align}
        \|\mathcal E_K(B)-B\|&\le2\|B\|.
        \label{eq:al_shell_identity_bound}
    \end{align}
\end{theorem}
\begin{proof}
    \leanok
    \uses{lem:al_site_expectation_order,def:al_local_root_channel}
    The expectation preserves the interval $[0,\mathbb 1]$.
    Thus the roots are self-adjoint and
    $G^2+R^2=\mathbb 1=G_K^2+R_K^2$.
    The map $V\psi=\widehat G\psi\oplus\widehat R\psi$ is an
    isometry and $\mathcal E(B)=V^\dagger(B\oplus B)V$, giving
    (\ref{eq:al_shell_full_contraction}). The same argument for
    $G_K,R_K$ gives~(\ref{eq:al_shell_local_contraction}). Finally,
    $\|\mathcal E_K(B)-B\|\le\|\mathcal E_K(B)\|+\|B\|$.
\end{proof}

\begin{theorem}[Channel error from localization of the constraint]
    \label{thm:al_root_channel_error}
    \uses{def:al_local_root_channel}
    \lean{TNLean.PEPS.AreaLaw.norm_spectatorRootChannel_sub_localRootChannel_le,
        TNLean.PEPS.AreaLaw.norm_localRootChannel_sub_localRootChannel_le}
    \leanok
    Suppose $\|k-E_K(k)\|\le\varepsilon_K$ for a real number
    $\varepsilon_K$. Then
    \begin{align}
        \|\mathcal E(B)-\mathcal E_K(B)\|
        &\le4\sqrt{\varepsilon_K}\,\|B\|.
        \label{eq:al_shell_localization_error}
    \end{align}
    If also $\|k-E_L(k)\|\le\varepsilon_L$, then
    \begin{align}
        \|\mathcal E_K(B)-\mathcal E_L(B)\|
        &\le4\bigl(\sqrt{\varepsilon_K}+\sqrt{\varepsilon_L}\bigr)\|B\|.
        \label{eq:al_shell_two_region_error}
    \end{align}
    The sets $K,L$ need not be nested. Both bounds are uniform over
    all finite spectator spaces.
\end{theorem}
\begin{proof}
    \leanok
    \uses{lem:al_site_expectation_order,def:al_local_root_channel}
    All four roots in~(\ref{eq:al_shell_full_roots}) and
    (\ref{eq:al_shell_local_roots}) have norm at most one.
    The square-root H\"older inequality gives
    $\|G-G_K\|\le\sqrt{\varepsilon_K}$ and
    $\|R-R_K\|\le\sqrt{\varepsilon_K}$.
    For physical operators $A,A'$, expand
    \begin{align}
        \widehat A B\widehat A-\widehat A'B\widehat A'
        &=(\widehat A-\widehat A')B\widehat A
            +\widehat A'B(\widehat A-\widehat A'),
        \label{eq:al_shell_root_difference}
    \end{align}
    \begin{align}
        \|\widehat A B\widehat A-\widehat A'B\widehat A'\|
        &\le(\|A\|+\|A'\|)\|A-A'\|\,\|B\|.
        \label{eq:al_shell_root_bound}
    \end{align}
    Apply~(\ref{eq:al_shell_root_bound}) to $(G,G_K)$ and $(R,R_K)$
    to obtain~(\ref{eq:al_shell_localization_error}). The triangle
    inequality through $\mathcal E(B)$ gives
    (\ref{eq:al_shell_two_region_error}). The hypotheses themselves
    imply $\varepsilon_K,\varepsilon_L\ge0$.
\end{proof}

\begin{theorem}[Fixing the regional commutant and intertwining commutators]
    \label{thm:al_root_channel_commutant}
    \uses{def:al_local_root_channel}
    \lean{TNLean.PEPS.AreaLaw.localRootChannel_eq_self_of_forall_commute,
        TNLean.PEPS.AreaLaw.localRootChannel_commutator_of_forall_commute,
        TNLean.PEPS.AreaLaw.localRootChannel_commutator}
    \leanok
    If $B$ commutes with every $\widehat C$ for $C\in\mathcal A_K$,
    then $\mathcal E_K(B)=B$. If an operator $A$ on
    $\mathcal H\otimes\mathcal K$ commutes with every such
    $\widehat C$, then
    \begin{align}
        [A,\mathcal E_K(B)]&=\mathcal E_K([A,B]).
        \label{eq:al_shell_commutator}
    \end{align}
    In particular, every physical $U\in\mathcal A_{I\setminus K}$
    satisfies $[\widehat U,\mathcal E_K(B)]
    =\mathcal E_K([\widehat U,B])$.
\end{theorem}
\begin{proof}
    \leanok
    \uses{lem:al_site_expectation_order,def:al_local_root_channel}
    The operator $k_K$ belongs to $\mathcal A_K$. Functional calculus
    preserves commutation with the complementary algebra, so both
    $G_K,R_K$ belong to $\mathcal A_K$ as well. Under the first
    hypothesis,
    \begin{align}
        \mathcal E_K(B)
        &=B(\widehat G_K^{\,2}+\widehat R_K^{\,2})=B.
        \label{eq:al_shell_commutant_fixed}
    \end{align}
    Under the second hypothesis, move $A$ past each root to obtain
    \begin{align}
        [A,\widehat G_K B\widehat G_K]
        &=\widehat G_K[A,B]\widehat G_K. \notag
    \end{align}
    The same identity holds for $R_K$; their sum proves
    (\ref{eq:al_shell_commutator}). Operators supported on disjoint
    physical regions commute, which gives the final assertion.
\end{proof}

\begin{theorem}[Oscillation outside the region]
    \label{thm:al_root_channel_outside_oscillation}
    \uses{def:al_local_root_channel,def:qc_site_oscillation}
    \lean{TNLean.PEPS.AreaLaw.siteOscillation_localRootChannel_le_of_notMem}
    \leanok
    If $y\notin K$, then $d_y(\mathcal E_K(B))\le d_y(B)$.
\end{theorem}
\begin{proof}
    \leanok
    \uses{thm:al_root_channel_commutant,thm:al_root_channel_contraction,
        thm:qc_site_oscillation_bounds}
    For every unitary $U\in\mathcal A_{\{y\}}$,
    \begin{align}
        \|[\widehat U,\mathcal E_K(B)]\|
        &=\|\mathcal E_K([\widehat U,B])\|
        \le\|[\widehat U,B]\|\le d_y(B).
        \label{eq:al_shell_outside_commutator}
    \end{align}
    Taking the supremum over $U$ proves the assertion, as in
    \cite[Section~10]{OpenAI2026AreaLaw}.
\end{proof}

\subsection{Differences of successive local channels}

\begin{definition}[Local channel shells]
    \label{def:al_local_channel_shell}
    \uses{def:al_local_root_channel}
    \lean{TNLean.PEPS.AreaLaw.localRootChannelShell,
        TNLean.PEPS.AreaLaw.localRootChannelShell_sub}
    \leanok
    Let $(K_l)_{l\ge0}$ be any sequence of physical regions and put
    $\mathcal E_l=\mathcal E_{K_l}$. Define
    \begin{align}
        \Delta_0(B)&=\mathcal E_0(B)-B,
        \label{eq:al_shell_zero}
    \end{align}
    \begin{align}
        \Delta_{l+1}(B)&=\mathcal E_{l+1}(B)-\mathcal E_l(B).
        \label{eq:al_shell_successor}
    \end{align}
    Equivalently, take $\mathcal E_{-1}=\mathrm{id}$ and
    $\Delta_l=\mathcal E_l-\mathcal E_{l-1}$, as in
    \cite[Section~10]{OpenAI2026AreaLaw}.
    Distributivity gives $\Delta_l(B-C)=\Delta_l(B)-\Delta_l(C)$.
\end{definition}

\begin{theorem}[Norm bounds for the channel shells]
    \label{thm:al_channel_shell_norm}
    \uses{def:al_local_channel_shell}
    \lean{TNLean.PEPS.AreaLaw.norm_localRootChannelShell_zero_le,
        TNLean.PEPS.AreaLaw.norm_localRootChannelShell_succ_le}
    \leanok
    For every sequence of regions, $\|\Delta_0(B)\|\le2\|B\|$.
    Fix $l\ge0$ and real numbers $\varepsilon_{\rm next},
    \varepsilon_{\rm prev}$ such that
    \begin{align}
        \|k-E_{K_{l+1}}(k)\|&\le\varepsilon_{\rm next},
        \label{eq:al_shell_next_error}
    \end{align}
    \begin{align}
        \|k-E_{K_l}(k)\|&\le\varepsilon_{\rm prev}.
        \label{eq:al_shell_previous_error}
    \end{align}
    Then
    \begin{align}
        \|\Delta_{l+1}(B)\|
        &\le4\bigl(\sqrt{\varepsilon_{\rm next}}
            +\sqrt{\varepsilon_{\rm prev}}\bigr)\|B\|.
        \label{eq:al_shell_successor_bound}
    \end{align}
\end{theorem}
\begin{proof}
    \leanok
    \uses{thm:al_root_channel_contraction,thm:al_root_channel_error}
    Apply~(\ref{eq:al_shell_identity_bound}) to $K_0$ for the
    zeroth shell and~(\ref{eq:al_shell_two_region_error}) to
    $K_{l+1},K_l$ for the successor shell.
\end{proof}

\begin{theorem}[Shells vanish on the physical regional commutant]
    \label{thm:al_channel_shell_commutant}
    \uses{def:al_local_channel_shell}
    \lean{TNLean.PEPS.AreaLaw.localRootChannelShell_eq_zero_of_forall_commute}
    \leanok
    Suppose the regions are increasing: $K_l\subseteq K_m$ whenever
    $l\le m$. If $B$ commutes with $\widehat A$ for every
    $A\in\mathcal A_{K_l}$, then $\Delta_l(B)=0$.
\end{theorem}
\begin{proof}
    \leanok
    \uses{thm:al_root_channel_commutant,def:al_local_channel_shell}
    For $l=0$, the commutant-fixing identity gives
    $\Delta_0(B)=B-B=0$. For $l=m+1$, the inclusion
    $\mathcal A_{K_m}\subseteq\mathcal A_{K_{m+1}}$ implies
    that both $\mathcal E_m$ and $\mathcal E_{m+1}$ fix $B$.
    Their difference therefore vanishes.
\end{proof}

\begin{theorem}[Physical-site oscillation bound for channel shells]
    \label{thm:al_channel_shell_oscillation}
    \uses{def:al_local_channel_shell,def:qc_site_oscillation}
    \lean{TNLean.PEPS.AreaLaw.norm_localRootChannelShell_zero_le_sum_siteOscillation,
        TNLean.PEPS.AreaLaw.norm_localRootChannelShell_succ_le_sum_siteOscillation}
    \leanok
    Suppose $(K_l)_{l\ge0}$ is increasing. Then
    \begin{align}
        \|\Delta_0(B)\|&\le2\sum_{y\in K_0}d_y(B).
        \label{eq:al_shell_zero_oscillation}
    \end{align}
    Fix $l\ge0$ and real numbers $\varepsilon_{\rm next},
    \varepsilon_{\rm prev}$ satisfying~(\ref{eq:al_shell_next_error})
    and~(\ref{eq:al_shell_previous_error}). Then
    \begin{align}
        \|\Delta_{l+1}(B)\|
        &\le4\bigl(\sqrt{\varepsilon_{\rm next}}
            +\sqrt{\varepsilon_{\rm prev}}\bigr)
            \sum_{y\in K_{l+1}}d_y(B).
        \label{eq:al_shell_successor_oscillation}
    \end{align}
    These bounds hold for every finite spectator and every complex
    operator $B$, without a Hermiticity assumption.
\end{theorem}
\begin{proof}
    \leanok
    \uses{thm:al_channel_shell_norm,thm:al_channel_shell_commutant,
        def:al_local_channel_shell,thm:qc_spectator_commutant,
        thm:qc_spectator_error,thm:qc_site_oscillation_bounds}
    Fix a shell index $m$ and a nonnegative coefficient $\eta$ from
    Theorem~\ref{thm:al_channel_shell_norm}, so
    $\|\Delta_m(C)\|\le\eta\|C\|$ for every $C$.
    Average precisely the physical sites in $K_m$, leaving the
    spectator untouched, by setting
    $\overline B=E_{I\setminus K_m}^{\mathcal K}(B)$.
    Theorems~\ref{thm:qc_spectator_commutant} and
    \ref{thm:qc_spectator_error} give, for every
    $A\in\mathcal A_{K_m}$,
    \begin{align}
        [\overline B,\widehat A]&=0,
        \label{eq:al_shell_averaged_commutant}
    \end{align}
    \begin{align}
        \|B-\overline B\|&\le\sum_{y\in K_m}d_y(B).
        \label{eq:al_shell_averaging_error}
    \end{align}
    Hence $\Delta_m(\overline B)=0$, and subtraction yields
    \begin{align}
        \|\Delta_m(B)\|
        &=\|\Delta_m(B-\overline B)\|
        \le\eta\|B-\overline B\|
        \le\eta\sum_{y\in K_m}d_y(B).
        \label{eq:al_shell_oscillation_bound}
    \end{align}
    For $m=0$ use $\eta=2$. For $m=l+1$ use
    $\eta=4(\sqrt{\varepsilon_{\rm next}}
    +\sqrt{\varepsilon_{\rm prev}})$, which is nonnegative.
    This is the commutant-averaging argument for the channel-shell
    estimate in \cite[Section~10]{OpenAI2026AreaLaw}.
\end{proof}
